\subsection{仿紧空间}

定义 6. 拓扑空间称为仿紧的，如果它是豪斯多夫的且满足以下公理：

(PC) X 的每个开覆盖 \(\Re\) 都有一个局部有限开加细 \(\Re'\)。（《集合论》，第二章，§ 4，第 6 目，定义 5）。

显然每个紧空间都是仿紧的。每个离散空间 X 都是仿紧的，因为由 X 的所有单点集组成的开覆盖是局部有限的，且比 X 的每个开覆盖都细。

命题 16. 仿紧空间 X 的每个闭子空间 F 都是仿紧的。

当然 F 是豪斯多夫的。另一方面，若 \((\mathbf{V}_{\iota})\) 是子空间 F 中的一个开覆盖，则每个 \(V_{\iota}\) 形如 \(V_{\iota} = U_{\iota} \cap F\)，其中 \(U_{\iota}\) 在 X 中是开的。考虑由 \(\complement F\) 与诸 \(U_{\iota}\) 组成的 X 的开覆盖 R；既然 X 是仿紧的，R 有一个局部有限加细 \(R'\)，而属于 \(R'\) 的诸集合与 F 的交构成 F 的一个局部有限开覆盖，它比给定的覆盖 \((\mathbf{V}_{\iota})\) 更细。

另一方面，紧空间的开子空间未必是仿紧的（习题 11）。

命题 17. 仿紧空间与紧空间的积是仿紧的。

设 X 是仿紧空间，Y 是紧空间，R 是 \(X \times Y\) 的一个开覆盖。对每个 \((x, y) \in X \times Y\)，存在 x 在 X 中的一个开邻域 \(V(x, y)\) 与 y 在 Y 中的一个开邻域 \(W(x, y)\)，使得 \(V(x, y) \times W(x, y)\) 包含于属于 R 的某个集合中。对每个 \(x \in X\)，当 y 跑遍 Y 时诸集合 \(W(x, y)\) 构成 Y 的一个开覆盖；因此存在有限多个点

\[
y _ {i} \in Y \quad (1 \leqslant i \leqslant n (x))
\]

使得诸 \(\mathbf{W}(x, y_i)\) 覆盖 \(\mathbf{Y}\)。设 \(\mathbf{U}(x)\) 表示

\[
\bigcap_ {i = 1} ^ {n (x)} \mathrm{V} (x, y _ {i});
\]

则诸开集 \(\mathbf{U}(x) \times \mathbf{W}(x, y_i)\) 中的每一个都包含于 \(\Re\) 的某个集合中。现在设 \((\mathbf{T}_{\iota})_{\iota \in \mathbf{I}}\) 是 \(\mathbf{X}\) 的一个局部有限开覆盖，它加细覆盖 \((\mathbf{U}(x))_{x \in \mathbf{X}}\)。对每个 \(\iota \in \mathbf{I}\)，设 \(x_{\iota}\) 是 \(\mathbf{X}\) 中一点，使得 \(\mathbf{T}_{\iota} \subset \mathbf{U}(x_{\iota})\)，并用 \(\mathbf{S}_{\iota,k}\) 表示诸集合 \(\mathbf{W}(x_{\iota}, y_k)\)（它们对应于这个点 \(x_{\iota}\)）（\(1 \leqslant k \leqslant n(x_{\iota})\)）。显然诸集合

\[
\mathrm{T} _ {\iota} \times \mathrm{S} _ {\iota, k} \quad (\iota \in \mathrm{I}, 1 \leqslant k \leqslant n (x _ {\iota}) \text {   for   each   } \iota \in \mathrm{I})
\]

构成 \(\mathbf{X} \times \mathbf{Y}\) 的一个加细 \(\Re\) 的开覆盖，而只要证明这个覆盖是局部有限的，证明即可完成。设 \((x, y)\) 是 \(\mathbf{X} \times \mathbf{Y}\) 的任意一点；存在邻域 \(\mathbf{Q}\)（即 \(x\) 的邻域），它只与有限多个集合 \(\mathbf{T}_{\iota}\) 相交，因而邻域 \(\mathbf{Q} \times \mathbf{Y}\)（\((x, y)\) 的）只与有限多个集合 \(\mathbf{T}_{\iota} \times \mathbf{S}_{\iota,k}\) 相交。

另一方面，两个仿紧空间的积未必是仿紧的（见第九章，§ 5，习题 16）。

命题 18. 仿紧空间族 \((\mathbf{X}_{\iota})_{\iota\in I}\) 的和（§ 2，第 4 目，例 3）是仿紧的。

设 X 是诸 \(X_{\iota}\) 的和，\((\mathrm{V}_{\lambda})_{\lambda\in\mathbf{L}}\) 是 X 的一个开覆盖。由开集 \(X_{\iota}\cap V_{\lambda}\) 组成的覆盖比 \((\mathrm{V}_{\lambda})\) 更细。对每个 \(\iota\in I\)，设 \((\mathrm{U}_{\iota,\mu})_{\mu\in\mathbf{M}_{\iota}}\) 是覆盖 \((\mathrm{V}_{\lambda}\cap\mathrm{X}_{\iota})_{\lambda\in\mathbf{L}}\) 的一个局部有限开加细；则由

\[
\mathbf {U} _ {\iota , \mu} \quad (\iota \in \mathbf {I}, \mu \in \mathbf {M} _ {\iota} \text {   for   each   } \iota \in \mathbf {I})
\]

组成的 X 的开覆盖是局部有限的，且加细原来的覆盖 \((\mathbf{V}_{\lambda})\)。

定理 5. 局部紧空间 X 是仿紧的，当且仅当 X 是局部紧 σ 紧空间族的一个和。

设 X 是仿紧的，并对每个 \(x \in X\) 设 \(V_x\) 是 x 在 X 中的一个相对紧开邻域。则由假设，存在 X 的一个局部有限开覆盖 \((\mathrm{U}(\alpha))_{\alpha \in \Lambda}\)，它加细覆盖 \((\mathrm{V}_x)_{x \in \mathbf{X}}\)。因此诸 \(\mathrm{U}(\alpha)\) 都是相对紧的。X 的每个紧子集 K 只与有限多个集合 \(\mathrm{U}(\alpha)\) 相交，因为非空集合 \(\mathrm{U}(\alpha) \cap \mathrm{K}\) 构成紧空间 X 的一个局部有限开覆盖；故它们必定只有有限个（第 1 目）。现在设 R 是 X 的两点 x, y 之间的如下关系：“存在 A 中的有限指标序列 \((\alpha_i)_{1 \leqslant i \leqslant n}\)，使得 \(x \in \mathrm{U}(\alpha_1)\)，\(y \in \mathrm{U}(\alpha_n)\)，且 \(\mathrm{U}(\alpha_i)\) 与 \(\mathrm{U}(\alpha_{i+1})\)（对 \(1 \leqslant i \leqslant n - 1\)）相交”。立即可以验证 R 是一个等价关系，且每个模 R 等价类都是 X 的开子集{[}因为诸 \(\mathrm{U}(\alpha)\) 是开集{]}。因此 X 是由模 R 的诸等价类构成的局部紧子空间（第 7 目，命题 13）的和，剩下只需证明这些子空间中的每一个都是族 \((\mathrm{U}(\alpha))\) 的一个可数子族之并。

设 \(x\) 是 \(\mathbf{X}\) 的任意一点，定义相对紧开子集的一个序列 \((\mathbf{C}_n)\)（\(\mathbf{X}\) 中的；对 \(n\) 用归纳法）如下：\(\mathbf{C}_1\) 是诸集合 \(\mathbf{U}(\alpha)\)（包含 \(x\) 的）之并，而对每个 \(n > 1\)，\(\mathbf{C}_n\) 是诸集合 \(\mathbf{U}(\alpha)\)（与 \(\mathbf{C}_{n-1}\) 相交的）之并。对 \(n\) 用归纳法立即可以验证，每个 \(\mathbf{C}_n\) 都是相对紧的，且是有限多个集合 \(\mathbf{U}(\alpha)\) 之并。此外，\(x\) 关于 \(\mathbf{R}\) 的等价类是诸 \(\mathbf{C}_n\) 之并：因为若 \((\alpha_i)_{1 \leqslant i \leqslant n}\) 是一个指标序列，使得 \(x \in \mathbf{U}(\alpha_1)\)，且 \(\mathbf{U}(\alpha_i)\) 与 \(\mathbf{U}(\alpha_{i+1})\)（对 \(1 \leqslant i \leqslant n - 1\)）相交，则对 \(i\) 用归纳法可以看出，\(\mathbf{U}(\alpha_i) \subset \mathbf{C}_i\)（对 \(1 \leqslant i \leqslant n\)）。由此推出模 \(\mathbf{R}\) 的诸等价类都是 \(\sigma\) 紧的，这就完成了定理第一部分的证明。

为证其逆，我们可以假设（由命题 18）X 是 σ 紧的。设 \(\Re = (G_{\lambda})_{\lambda \in L}\) 是 X 的任意一个开覆盖，\((U_n)\) 是 X 中具有第 9 目命题 15 所述性质的相对紧开集序列。设 \(K_n\) 表示紧集 \(\overline{U}_n - U_{n-1}(U_n = \emptyset\)，若 \(n \leqslant 0)\)。由构造，开集 \(U_{n+1} - \overline{U}_{n-2}\) 是 \(K_n\) 的一个邻域；因此对每个 \(x \in K_n\)，存在邻域 \(W_x\)（即 \(x\) 的邻域），它包含于某个集合 \(G_{\lambda}\) 中，且也包含于 \(U_{n+1} - \overline{U}_{n-2}\) 中。既然 \(K_n\) 是紧的，有限多个集合 \(W_x\) 便覆盖 \(K_n\)；设 \(H_{ni}(1 \leqslant i \leqslant p_n)\) 是这些集合。于是族 \(\Re'\)（由诸集合 \(H_{ni}(n \geqslant 1, 1 \leqslant i \leqslant p_n\)，对每个 \(n\)) 组成）是 X 的一个加细 \(\Re\) 的开覆盖，因此为完成证明，只须证明 \(\Re'\) 是局部有限的。设 \(z\) 是 X 的任意一点，\(n\) 是使 \(z \in U_n\) 的最小整数；则由于 \(z \notin U_{n-1}\)，存在 \(z\) 的一个邻域 T，它包含于 \(U_n\) 中且不与 \(U_{n-2}\) 相交。由此推出 T 只与那些集合 \(H_{mi}\)（对它们 \(n - 2 \leqslant m \leqslant n + 1\)）相交，即 T 只与 \(\Re'\) 的有限多个集合相交。

在证明过程中我们还建立了如下结果：

推论. 设 X 是局部紧仿紧空间。则 X 的每个开覆盖 R 都有一个由相对紧集组成的局部有限开加细 R'。若 X 是 σ 紧的，则可取 R' 为可数的。

